\begin{ExerciseBlock}
\ExerciseHeading{01}{Interpretation des Auslenkungs-Zeit-Gesetzes}{free_oscillations_displacement_law}
Ein Körper schwingt in $x$-Richtung. Seine Auslenkung folgt dem Zeitgesetz
\[
 x(t)=\qty{5}{\metre}\cos\left(3{,}6\pi\,\unit{\per\second}\,t+\frac{\pi}{3}\right).
\]
\begin{enumerate}
\item Geben Sie $x_0$, $\omega_0$, $f_0$, $T_0$ und den Nullphasenwinkel $\varphi_0$ an.
\item Berechnen Sie $x(t_0)$, $v(t_0)$ und $a(t_0)$ zum Zeitpunkt $t_0=\qty{2}{\second}$.
\end{enumerate}
\end{ExerciseBlock}
\UseExerciseSolution{free_oscillations_displacement_law}
